
\section{Language grammar (EBNF)}

We use EBNF\cite{ISO14977} to define the grammar of our functional language.
\begin{lstlisting}
(* Program *)
program = { declaration }, "render", expression

(* Declarations *)
declaration = "let", identifier, [":", type], "=", expression
            | "event", identifier, ":", type

(* Expressions *)
expression = literal | identifier | lambda | expression expression
            | "let", identifier, [":", type], "=", expression, "in", expression
            | "if" expression, "then", expression, "else", expression
            | expression binop expression
            | "[", expression, {";", expression}, "]"

lambda = "fn", "(", identifier, ":", type, ")", "->", expression

binop = "+" | "-" | "*" | "==" | "<" | "^"

literal = integer | "true" | "false" | string | "()"

(* Types *)
type = "Int" | "Bool" | "String" | "UI" | "Unit" | "List", type
      | "Behaviour", type | "Event", type | type "->" type | "(" type ")"

(* Identifiers and Primitive Values *)
identifier = ? [a-zA-Z0-9]+ ? ;

integer = ? [0-9]+ ? ;

string = '"', { string_char | escape }, '"'  ;

string_char = ? any Unicode character except " and \ ?  ;

escape = "\", ( '"' | "\" | "n" | "t" )  ;
\end{lstlisting}

\medskip
\noindent The langage includes identifiers for functions that produce UI elements and combinators for treating signal types. We list these below with their type.

\begin{center}
  \begin{tabular}{c | c}
    Name & Type\\
    \hline
    \texttt{show} & \texttt{Int -> String}\\
    \hline
    \texttt{text} & \texttt{String -> UI}\\
    \hline
    \texttt{button} & \texttt{String -> Event Unit -> UI}\\
    \hline
    \texttt{column}, \texttt{row} & \texttt{List UI -> UI}\\
    \hline
    \texttt{dyn} & \texttt{Behaviour UI -> UI}\\
    \hline
    \texttt{map} & \texttt{('a -> 'b) -> Behaviour 'a -> Behaviour 'b}\\
    \hline
    \texttt{map2} & \texttt{('a -> 'b -> 'c) -> Behaviour 'a -> Behaviour 'b -> Behaviour 'c}\\
    \hline
    \texttt{mapE} & \texttt{('a -> 'b) -> Event 'a -> Event 'b}\\
    \hline
    \texttt{filterE} & \texttt{('a -> Bool) -> Event 'a -> Event 'a}\\
    \hline
    \texttt{mergeWith} & \texttt{('a -> 'a -> 'a) -> Event 'a -> Event 'a -> Event 'a}\\
    \hline
    \texttt{fold} & \texttt{'b -> ('b -> 'a -> 'b) -> Event 'a -> Behaviour 'b}\\
    \hline
    \texttt{hold} & \texttt{'a -> Event 'a -> Behaviour 'a}\\
    \hline
    \texttt{snapshot} & \texttt{('a -> 'b -> 'c) -> Event 'a -> Behaviour 'b -> Event 'c}\\
\end{tabular}
\end{center}

The meta-variables \texttt{'a} are used as a placeholder to iterate over all possible types. In the langauge there is no polymorphism.
